\documentclass[11pt,a4paper]{article}
\usepackage[margin=2.2cm]{geometry}
\usepackage[utf8]{inputenc}
\usepackage[T1]{fontenc}
\usepackage{amsmath,amssymb,amsfonts,bm}
\usepackage{braket}
\usepackage{enumitem}
\usepackage{xcolor}
\usepackage{tikz}
\usetikzlibrary{quantikz2}

\newcommand{\pts}[1]{\textbf{(#1)}}
\newcommand{\tr}{\operatorname{tr}}
\newcommand{\id}{\mathbb{I}}
\newcommand{\sol}{\color{blue!60!black}}

\setlist[enumerate,1]{label=\alph*)}
\setlength{\parindent}{0pt}
\setlength{\parskip}{4pt}

\title{Quantum Technologies -- First partial exam\\[4pt] \large Practice exam B}
\date{}

\begin{document}
\maketitle
\vspace{-1.5cm}
\begin{center}
Duration: 2 hours. Total: 10 points.
\end{center}

%=====================================================================
\section{Open question (1pt)}

What does it mean that a set of quantum gates is \emph{universal}? Give an example of a universal
set of gates and explain why the Clifford group, generated by $\{H,S,\text{CNOT}\}$, is not universal.
What are the consequences of this for quantum computation? Finally, explain why the
Solovay--Kitaev theorem is relevant for building quantum computers in practice.
(Maximum 200 words, $\sim$1 page)

%=====================================================================
\section{Short questions (3pts)}

\begin{enumerate}
\item Consider the following circuit, where the dashed box is the oracle $U_f$ of the
Bernstein--Vazirani algorithm. Which function $f(\vec x)$ does the oracle implement? What is the
result of measuring the first three qubits, and with which probability? How many evaluations of $f$
would be needed classically to find the hidden string? \pts{1 p}
\begin{center}
\begin{quantikz}[row sep=0.25cm, column sep=0.35cm]
\lstick{$\ket{0}$} & & \gate{H} & \ctrl{3}\gategroup[4,steps=3,style={dashed,rounded corners,inner xsep=2pt},label style={label position=above}]{$U_f$} & & & \gate{H} & \meter{} \\
\lstick{$\ket{0}$} & & \gate{H} & & & & \gate{H} & \meter{} \\
\lstick{$\ket{0}$} & & \gate{H} & & \ctrl{1} & & \gate{H} & \meter{} \\
\lstick{$\ket{0}$} & \gate{X} & \gate{H} & \targ{} & \targ{} & \gate{X} & &
\end{quantikz}
\end{center}

\item Consider the following phase-kickback circuit, where $T=\begin{pmatrix}1&0\\0&e^{i\pi/4}\end{pmatrix}$.
\begin{center}
\begin{quantikz}[row sep=0.35cm, column sep=0.4cm]
\lstick{$\ket{0}$} & \gate{H} & \ctrl{1} & \gate{H} & \meter{} \\
\lstick{$\ket{\psi}$} & & \gate{T} & &
\end{quantikz}
\end{center}
What is the probability of measuring the first qubit in state $\ket0$ if the second qubit is in the
state $\ket\psi=\ket0$? And if it is in $\ket\psi=\ket1$? We know that $\ket\psi$ is either $\ket0$
or $\ket1$. How many times do we need to repeat the circuit to identify $\ket\psi$ with a
probability of error smaller than $5\%$? \pts{1 p}

\item We want to find a prime factor of $M=33$ using the Shor algorithm with $x=5$. Solving the
order-finding problem with a quantum computer, with $n=8$ qubits in the first register, we
measure $\vec y_0=\{10110011\}$. Use the continued fraction expansion to find the order of $x$
modulo $M$, and use it to find the prime factors of $M$. \pts{1 p}
\end{enumerate}

%=====================================================================
\section*{Problems}

\section{Phase-flip channel (3pts)}

A qubit is sent through a noisy channel described by the measurement operators
\[
M_0=\sqrt{p}\,\id,\qquad M_1=\sqrt{1-p}\,\sigma_z,\qquad 0\le p\le1 .
\]
\begin{enumerate}
\item Show that $\{M_0,M_1\}$ define a generalized measurement and find the associated POVM
elements $E_0,E_1$ \pts{0.5 p}.
\item For a general state $\rho$, calculate the probability of each outcome and the state after
the measurement for each outcome. What information about $\rho$ can we learn from this measurement?
\pts{0.5 p}
\item The outcome of the measurement is not recorded. Write the output state $\rho'$ and show that,
if $\rho=\frac12(\id+\bm r\cdot\bm\sigma)$, the Bloch vector of the output is
$\bm r'=\left((2p-1)r_x,\,(2p-1)r_y,\,r_z\right)$ \pts{0.5 p}.
\item Consider now the pure input state $\ket\psi=\cos\frac\theta2\ket0+e^{i\phi}\sin\frac\theta2\ket1$.
Calculate the purity of the output $\rho'$. For which input states, and for which values of $p$,
does the output remain pure? \pts{0.5 p}
\item Calculate the fidelity $F(\ket\psi\!\bra\psi,\rho')$ between the input and the output states \pts{0.5 p}.
\item Calculate the trace distance $d(\ket\psi\!\bra\psi,\rho')$. For $p=1/2$, which input states
are not affected by the channel, and which ones are mapped to the maximally mixed state? \pts{0.5 p}
\end{enumerate}

\section{Measuring half of an entangled state (3pts)}

Alice (qubit $A$, first) and Bob (qubit $B$, second) share the two-qubit state
\[
\ket\psi=\frac{1}{2\sqrt2}\left[\sqrt3\ket{00}+\sqrt3\ket{01}+\ket{10}-\ket{11}\right].
\]
\begin{enumerate}
\item Check that $\ket\psi$ is normalized and write its density matrix $\rho$ in the computational
basis. Is it a pure state? \pts{0.5 p}
\item Calculate the reduced density matrices $\rho_A$ and $\rho_B$, locate them in the Bloch sphere,
and calculate their purities \pts{0.5 p}.
\item Find the Schmidt decomposition of $\ket\psi$ and its entanglement entropy \pts{0.5 p}.
\item Alice measures $\sigma_x$ on her qubit. Calculate the probability of each outcome and the
state of Bob's qubit after the measurement, for each outcome \pts{0.5 p}.
\item If Alice does not tell Bob the result, what is the state of Bob's qubit after her measurement?
Compare it with the state before the measurement and discuss the result \pts{0.5 p}.
\item Bob does not know the result of Alice's measurement, but he wants to know which of the two
states of part d) he has. Can he distinguish them perfectly? Calculate their fidelity and the
maximum probability of success of an unambiguous discrimination strategy \pts{0.5 p}.
\end{enumerate}

%=====================================================================
\newpage
\begin{center}
\Large\textbf{Solutions}
\end{center}
\sol

\section*{1\quad Open question}

A set of gates $G$ is \textbf{universal} if any unitary $U$ can be approximated to arbitrary
precision $\epsilon$ by a circuit made only of gates from $G$. Any single-qubit unitary can be
decomposed as $U=e^{i\alpha}R_z(\beta)R_y(\gamma)R_z(\delta)$, and single-qubit gates together with
CNOT can build any $n$-qubit unitary. Therefore $\{R_z(\theta),R_y(\theta),\text{CNOT}\}$ is
universal. With a finite set, $\{H,T,\text{CNOT}\}$ (Clifford+T) is universal: $H$ and $T$ generate
rotations by an angle that is an irrational multiple of $2\pi$, which approximate any rotation.

The Clifford group $\{H,S,\text{CNOT}\}$ is not universal. It only maps Pauli operators to Pauli
operators, and, for example, $T=\sqrt S$ cannot be built from Clifford gates. The
\textbf{Gottesman--Knill theorem} states that Clifford circuits can be simulated efficiently on a
classical computer. So they cannot give a quantum advantage, and some non-Clifford resource
(``magic'', such as the $T$ gate) is needed.

Universality alone does not guarantee that the approximation is efficient. The
\textbf{Solovay--Kitaev theorem} states that, with a finite universal set, any single-qubit gate can
be approximated within $\epsilon$ using only $\mathcal{O}(\log^c(1/\epsilon))$ gates. Hardware can
therefore implement (and error-correct) a small fixed set of gates, at a modest, polylogarithmic cost.

\section*{2\quad Short questions}

\begin{enumerate}
\item The two CNOTs flip the ancilla when $x_0=1$ and when $x_2=1$, and the final $X$ flips it
unconditionally. So $U_f\ket{\vec x}\ket y=\ket{\vec x}\ket{y\oplus f(\vec x)}$ with
\[
f(\vec x)=x_0\oplus x_2\oplus1=\vec s\cdot\vec x\oplus1,\qquad \vec s=101 .
\]
With the ancilla in $\ket-$, the oracle acts as $\ket{\vec x}\ket-\mapsto(-1)^{f(\vec x)}\ket{\vec x}\ket-=-(-1)^{\vec s\cdot\vec x}\ket{\vec x}\ket-$:
\[
\ket{000}\ket1\xrightarrow{H^{\otimes4}}\frac1{\sqrt8}\sum_{\vec x}\ket{\vec x}\ket-
\xrightarrow{U_f}-\frac1{\sqrt8}\sum_{\vec x}(-1)^{\vec s\cdot\vec x}\ket{\vec x}\ket-
\xrightarrow{H^{\otimes3}}-\ket{101}\ket- .
\]
The constant term only adds an irrelevant global phase $-1$, so we measure $\vec s=\{101\}$
with probability 1, using a single evaluation of $f$. Classically one needs $n=3$ evaluations
($f(100),f(010),f(001)$) to read the bits of $\vec s$. Because of the unknown constant, one more
evaluation, $f(000)$, is needed, which makes 4 in total.

\item Following the phase-kickback analysis: if $\ket\psi$ is an eigenstate of $T$ with eigenvalue
$e^{i\vartheta}$,
\[
\ket0\ket\psi\xrightarrow{H}\tfrac1{\sqrt2}(\ket0+\ket1)\ket\psi\xrightarrow{CT}\tfrac1{\sqrt2}(\ket0+e^{i\vartheta}\ket1)\ket\psi
\xrightarrow{H}\tfrac12\left[(1+e^{i\vartheta})\ket0+(1-e^{i\vartheta})\ket1\right]\ket\psi,
\]
so $P(0)=\left|\frac{1+e^{i\vartheta}}2\right|^2=\cos^2\frac\vartheta2$.
\begin{itemize}
\item $\ket\psi=\ket0$: $\vartheta=0$, so $P(0)=1$.
\item $\ket\psi=\ket1$: $\vartheta=\pi/4$, so $P(0)=\cos^2\frac\pi8=\frac{2+\sqrt2}{4}\approx0.854$ and $P(1)\approx0.146$.
\end{itemize}
Strategy: if we obtain the outcome 1 at least once, $\ket\psi=\ket1$ with certainty. If all $m$
repetitions give 0, we guess $\ket0$. The only error happens when $\ket\psi=\ket1$ and all $m$ outcomes are 0:
\[
P_{\rm err}=\cos^{2m}\frac\pi8<0.05\ \Rightarrow\ m>\frac{\ln0.05}{\ln0.854}\approx18.9\ \Rightarrow\ m=19 .
\]

\item $\vec y_0=\{10110011\}$ corresponds to $y_0=128+32+16+2+1=179$, with $N=2^8=256$ and
$\frac1{2N}=\frac1{512}\approx0.00195$. Continued fraction expansion (Euclid's algorithm:
$256=1\cdot179+77$, $179=2\cdot77+25$, $77=3\cdot25+2$, $25=12\cdot2+1$):
\[
\frac{179}{256}=\cfrac{1}{1+\cfrac{1}{2+\cfrac{1}{3+\cfrac{1}{12+\cfrac12}}}}=[0;1,2,3,12,2],
\qquad\text{convergents: }0,\ 1,\ \frac23,\ \frac7{10},\ \frac{86}{123},\ \frac{179}{256}.
\]
$\left|\frac{179}{256}-\frac23\right|\approx0.033>\frac1{512}$ (and indeed $5^3=125\equiv26\not\equiv1 \bmod 33$), while
$\left|\frac{179}{256}-\frac7{10}\right|=\frac{2}{2560}\approx0.00078<\frac1{512}$. The candidate is $r=10<N$.
Check: $5^1,\dots,5^{10}\bmod 33=5,25,26,31,23,16,14,4,20,1$, so $r=10$ is the order of 5 modulo 33.

Shor's algorithm: $r$ is even and $x^{r/2}=5^5\equiv23\bmod33$, with $23\not\equiv-1\equiv32$. Then
$M_\pm=x^{r/2}\pm1$ give $M_+=24$ and $M_-=22$, and
\[
\gcd(24,33)=3,\qquad\gcd(22,33)=11\qquad\Rightarrow\qquad 33=3\times11 .
\]
(Euclid: $33=1\cdot24+9$, $24=2\cdot9+6$, $9=1\cdot6+3$, $6=2\cdot3$; and $33=1\cdot22+11$, $22=2\cdot11$.)
\end{enumerate}

\section*{3\quad Phase-flip channel}

\begin{enumerate}
\item Completeness: $M_0^\dagger M_0+M_1^\dagger M_1=p\,\id+(1-p)\sigma_z^2=p\,\id+(1-p)\id=\id$, since
$\sigma_z^\dagger=\sigma_z$ and $\sigma_z^2=\id$. So the operators define a generalized measurement. The
POVM elements are $E_m=M_m^\dagger M_m$:
\[
E_0=p\,\id,\qquad E_1=(1-p)\,\id,
\]
which are positive semi-definite ($0\le p\le1$) and sum to $\id$.

\item $p(0)=\tr[E_0\rho]=p\,\tr\rho=p$ and $p(1)=\tr[E_1\rho]=1-p$. Both are independent of $\rho$.
The post-measurement states are
\[
\rho_0'=\frac{M_0\rho M_0^\dagger}{p(0)}=\frac{p\,\rho}{p}=\rho,\qquad
\rho_1'=\frac{M_1\rho M_1^\dagger}{p(1)}=\sigma_z\rho\,\sigma_z .
\]
Since the outcome probabilities do not depend on $\rho$, the measurement gives \textbf{no
information} about the state. It only tells us whether a $\sigma_z$ (phase flip) was applied.

\item $\rho'=\sum_mM_m\rho M_m^\dagger=p\,\rho+(1-p)\,\sigma_z\rho\,\sigma_z$. Using
$\sigma_z\sigma_x\sigma_z=-\sigma_x$, $\sigma_z\sigma_y\sigma_z=-\sigma_y$ and $\sigma_z\sigma_z\sigma_z=\sigma_z$
(from $\sigma_i\sigma_j=\delta_{ij}\id+i\epsilon_{ijk}\sigma_k$):
\[
\sigma_z\rho\,\sigma_z=\tfrac12\left(\id-r_x\sigma_x-r_y\sigma_y+r_z\sigma_z\right),
\]
\[
\rho'=\tfrac12\left[\id+(p-(1-p))(r_x\sigma_x+r_y\sigma_y)+r_z\sigma_z\right]
\ \Rightarrow\ \bm r'=\left((2p-1)r_x,\,(2p-1)r_y,\,r_z\right).
\]
The Bloch ball is squeezed in the $xy$ plane by a factor $|2p-1|$, while the $z$ axis is left
unchanged: the channel destroys coherences, but not the populations.

\item For $\ket\psi$, $\bm r=(\sin\theta\cos\phi,\sin\theta\sin\phi,\cos\theta)$, so
$\bm r'=\left((2p-1)\sin\theta\cos\phi,\,(2p-1)\sin\theta\sin\phi,\,\cos\theta\right)$ and
\[
\gamma(\rho')=\frac{1+|\bm r'|^2}{2}=\frac{1+(2p-1)^2\sin^2\theta+\cos^2\theta}{2}=1-2p(1-p)\sin^2\theta .
\]
The output is pure ($\gamma=1$) iff $p(1-p)\sin^2\theta=0$. This happens for $p=1$ (no flip), for $p=0$
(a deterministic $\sigma_z$, which is unitary), or for $\sin\theta=0$, i.e.\ input states $\ket0$ or $\ket1$
(the poles, which are eigenstates of $\sigma_z$).

\item Since the input is pure, $F=\bra\psi\rho'\ket\psi=\frac{1+\bm r\cdot\bm r'}{2}$
(using $\tr[\rho\sigma]=\frac{1+\bm r\cdot\bm s}{2}$ for qubits):
\[
F=\frac{1+(2p-1)\sin^2\theta+\cos^2\theta}{2}=1-(1-p)\sin^2\theta .
\]
Alternatively: $\bra\psi\rho'\ket\psi=p+(1-p)|\bra\psi\sigma_z\ket\psi|^2=p+(1-p)\cos^2\theta$, which is the same.
It is 1 at the poles and minimal, $F=p$, on the equator ($\theta=\pi/2$).

\item For qubits, $d=\frac12|\bm r-\bm r'|$:
\[
\bm r-\bm r'=\left(2(1-p)\sin\theta\cos\phi,\ 2(1-p)\sin\theta\sin\phi,\ 0\right)
\ \Rightarrow\ d(\ket\psi\!\bra\psi,\rho')=(1-p)\sin\theta .
\]
For $p=1/2$: $\bm r'=(0,0,\cos\theta)$. The states $\ket0$ and $\ket1$ ($\theta=0,\pi$) are not
affected ($d=0$, $F=1$). The equatorial states ($\theta=\pi/2$, e.g.\ $\ket\pm$, $\ket{R}$, $\ket{L}$) are
mapped to $\bm r'=\bm 0$, i.e.\ to the maximally mixed state $\id/2$. For these, $d=1/2$ and $F=1/2$.
\end{enumerate}

\section*{4\quad Measuring half of an entangled state}

\begin{enumerate}
\item $\braket{\psi|\psi}=\frac18(3+3+1+1)=1$. In the basis $\{\ket{00},\ket{01},\ket{10},\ket{11}\}$,
$\ket\psi=\frac1{2\sqrt2}(\sqrt3,\sqrt3,1,-1)^T$ and
\[
\rho=\ket\psi\!\bra\psi=\frac18\begin{pmatrix}3&3&\sqrt3&-\sqrt3\\3&3&\sqrt3&-\sqrt3\\\sqrt3&\sqrt3&1&-1\\-\sqrt3&-\sqrt3&-1&1\end{pmatrix}.
\]
It is a pure state by construction, since $\rho=\ket\psi\!\bra\psi$ is a rank-1 projector. Equivalently, $\rho^2=\rho$ and $\tr\rho^2=1$.

\item Using $(\rho_A)_{ij}=\rho_{i0,j0}+\rho_{i1,j1}$ and $(\rho_B)_{ij}=\rho_{0i,0j}+\rho_{1i,1j}$:
\[
\rho_A=\frac18\begin{pmatrix}3+3&\sqrt3-\sqrt3\\\sqrt3-\sqrt3&1+1\end{pmatrix}=\begin{pmatrix}3/4&0\\0&1/4\end{pmatrix},
\qquad
\rho_B=\frac18\begin{pmatrix}3+1&3-1\\3-1&3+1\end{pmatrix}=\frac12\begin{pmatrix}1&1/2\\1/2&1\end{pmatrix}.
\]
Bloch vectors: $\bm r_A=(0,0,\tfrac12)$, on the $z$ axis halfway to $\ket0$, and
$\bm r_B=(\tfrac12,0,0)$, on the $x$ axis halfway to $\ket+$. Both are inside the Bloch ball. Purities:
$\gamma(\rho_A)=\gamma(\rho_B)=\frac{1+1/4}{2}=\frac58<1$. They are equal, as expected for a pure global state. Both
reduced states are mixed, so $\ket\psi$ is entangled.

\item $\rho_A$ is already diagonal: $\eta_1^2=\frac34$ with $\ket{1_A}=\ket0$, and $\eta_2^2=\frac14$ with
$\ket{2_A}=\ket1$. Projecting onto them:
\[
\bra0_A\ket\psi=\frac{\sqrt3}{2\sqrt2}(\ket0+\ket1)=\frac{\sqrt3}{2}\ket+,\qquad
\bra1_A\ket\psi=\frac{1}{2\sqrt2}(\ket0-\ket1)=\frac12\ket- .
\]
\[
\boxed{\ket\psi=\frac{\sqrt3}{2}\ket0\ket++\frac12\ket1\ket-}
\]
The Schmidt coefficients are $\eta_1=\frac{\sqrt3}{2}$ and $\eta_2=\frac12$, with bases
$\{\ket0,\ket1\}_A$ and $\{\ket+,\ket-\}_B$. As a check, $\rho_B=\frac34\ket+\!\bra++\frac14\ket-\!\bra-$
agrees with part b). The Schmidt rank is 2 (entangled), and
\[
S_{EE}=-\tfrac34\log_2\tfrac34-\tfrac14\log_2\tfrac14\approx0.81 .
\]

\item We write Alice's states in the $\sigma_x$ eigenbasis, $\ket0=\frac{\ket++\ket-}{\sqrt2}$ and $\ket1=\frac{\ket+-\ket-}{\sqrt2}$:
\[
\ket\psi=\frac1{\sqrt2}\ket+\otimes\left(\tfrac{\sqrt3}{2}\ket++\tfrac12\ket-\right)+\frac1{\sqrt2}\ket-\otimes\left(\tfrac{\sqrt3}{2}\ket+-\tfrac12\ket-\right).
\]
With projectors $P_\pm=\ket\pm\!\bra\pm\otimes\id$, the outcomes $\pm1$ have probabilities
$p(\pm1)=\bra\psi P_\pm\ket\psi=\frac12\left(\frac34+\frac14\right)=\frac12$, and Bob's qubit is left in
\[
\ket{\phi_\pm}=\frac{\sqrt3}{2}\ket+\pm\frac12\ket- .
\]
(In the computational basis, $\ket{\phi_+}=\frac{(\sqrt3+1)\ket0+(\sqrt3-1)\ket1}{2\sqrt2}=\cos\frac\pi{12}\ket0+\sin\frac\pi{12}\ket1$ and
$\ket{\phi_-}=\sin\frac\pi{12}\ket0+\cos\frac\pi{12}\ket1$.)

\item If the outcome is not communicated, Bob has the ensemble $\{(\frac12,\ket{\phi_+}),(\frac12,\ket{\phi_-})\}$:
\[
\rho_B'=\tfrac12\ket{\phi_+}\!\bra{\phi_+}+\tfrac12\ket{\phi_-}\!\bra{\phi_-}
=\tfrac34\ket+\!\bra++\tfrac14\ket-\!\bra-=\rho_B ,
\]
because the cross terms $\pm\frac{\sqrt3}4(\ket+\!\bra-+\ket-\!\bra+)$ cancel. Bob's state is unchanged:
nothing Alice does locally can change the statistics of Bob's measurements (\emph{no-signalling}).
Bob can only notice the correlations once Alice sends him her results through a classical channel.
In general, $\rho_B=\tr_A\rho$ does not depend on the measurement made on $A$.

\item Bob's two possible states are not orthogonal:
\[
\braket{\phi_+|\phi_-}=\frac34-\frac14=\frac12\neq0 .
\]
So no measurement can distinguish them with certainty. Their fidelity (pure states) is
$F=|\braket{\phi_+|\phi_-}|^2=\frac14$. With the unambiguous discrimination POVM
$E_1=a(\id-\ket{\phi_-}\!\bra{\phi_-})$, $E_2=a(\id-\ket{\phi_+}\!\bra{\phi_+})$, $E_3=\id-E_1-E_2$, the
largest allowed value is $a=\frac1{1+|\braket{\phi_+|\phi_-}|}=\frac23$. For equiprobable states this gives
\[
p_{\rm success}=1-|\braket{\phi_+|\phi_-}|=\frac12 .
\]
Half of the time Bob identifies his state with certainty. The other half he gets the inconclusive outcome $E_3$.
\end{enumerate}

\end{document}
